\begin{table}[htbp]
\raggedright
\caption{Profile moderators of attacked effectivity: controlled contrasts, SEM paths, and descriptive susceptibility weights.}
\label{tab:multivariate}
{
\footnotesize
\setlength{\tabcolsep}{4pt}
\renewcommand{\arraystretch}{1.08}
\resizebox{\linewidth}{!}{%
\begin{tabular}{p{0.34\linewidth}rrrrrrrrr}
\toprule
Moderator & Role & Controlled mean b & Controlled p & Ridge mean b & Weight \% & SEM mean b & SEM mean |b| & SEM min p & SEM paths \\
\midrule
Big Five Openness To Experience Mean \% & core & -1.834 & 0.028 &  &  & -1.894 & 2.574 & 0.074 & 10.000 \\
Big Five Neuroticism Mean \% & core & 1.076 & 0.162 &  &  & 1.482 & 3.032 & 0.027 & 10.000 \\
Sex Other & core & 1.748 & 0.425 & 0.033 & 1.023 & 4.729 & 7.383 & 0.004 & 10.000 \\
Big Five Conscientiousness Mean \% & core & 0.494 & 0.689 &  &  & 0.962 & 4.505 & 0.018 & 10.000 \\
Sex Female & core & 0.454 & 0.844 & -0.108 & 0.961 & 2.852 & 9.409 & 0.032 & 10.000 \\
Big Five Extraversion Mean \% & exploratory & -0.248 & 0.829 &  &  &  &  &  &  \\
Big Five Agreeableness Mean \% & exploratory & -0.307 & 0.786 &  &  &  &  &  &  \\
\bottomrule
\end{tabular}
}
}
\par\smallskip
{\normalsize Note. Controlled coefficients summarize moderator associations with mean attacked effectivity. SEM columns summarize the repeated-outcome path model across the attacked opinion leaves. Weight percentages come from the target-conditional regularized aggregation used to compute the post hoc susceptibility index.}
\end{table}